\subsection*{六次Prym像の再利用可能な標準形}

階数七の係数Gramの核は二次元pencilである。
行列式が相異なる三根を持つ開部分では、これを対角の跡零pencilに正規化できる。
正規化後の $P$ と三次元の一次直交syzygy基底 $R$ は
論文Vの「九本の余因子方程式」で与えた固定行列となり、
全ての変形情報は十九変数の対称形式 $A$ に移る。
負ブロックは $B=RA^{-1}$ であり、残る条件は
\[
 \begin{aligned}
 (A^{-1})_{14}&=(A^{-1})_{15}=(A^{-1})_{16}=0,\\
 (A^{-1})_{23}&=(A^{-1})_{24}=(A^{-1})_{26}=0,\\
 (A^{-1})_{35}&=(A^{-1})_{36}=(A^{-1})_{45}=0.
 \end{aligned}
\]
である。添字は $0,\ldots,6$。
従って、後続の計算を百五十三変数・七十二方程式のチャートから、
十九変数・九余因子方程式のチャートへ縮小できる。
この正規化に必要な三根の順序付けは、最後に $\mathfrak S_3$ で除く。

$A_0$ は $(0)$ の成分が $1$ で、$(1,3),(2,5),(4,6)$ を交換する行列である。
$\det A_0=-1$。
上記の順序の九余因子の $(p,m,i,r,o,l,b,g,f)$ に関するJacobi行列は $I_9$ である。
従って局所計算には $(s-1,a,c,d,e,h,j,k,n,q)$ の十変数を自由変数として用いられる。
これは九個の変数を形式的・エタール局所的に消去できるという主張であり、
大域的な有理パラメータ表示を与えるものではない。
また $A_0$ ではトーラスの安定化群が一般点より大きいため、
十次元チャートの滑らかさだけから粗商の基点の滑らかさを推論しない。

正則核pencilを持つ階数七Gram軌跡は、十四次元のpencil軌道上の
十九次元線形空間族の開部分であり、滑らかな三十三次元空間である。
九方程式は既存成分に沿う一般点で余次元九を持つ。
選択した十次元成分の有効安定化群は
$\mathbb G_m^3\rtimes\mathfrak S_3$ であり、その商は七次元である。
全余因子軌跡の既約性や他の成分は未判定。
成分の選択は既存の整数基点 $B_-$ と同定した $A_0$ により固定する。

\paragraph{標準的な直線とPrym-canonical切断.}
核pencilの三つの特異元から、両三次元空間の固有直線分解を取り戻せる。
対応する三つの階数一テンソルが張る三次元空間 $E$ は根の順序に依存せず、
$\ell=E/K\subset (U\otimes V)/K$ は標準的な直線となる。
これは標準形の $\langle e_0\rangle$ である。
$v\in\ell\setminus\{0\}$ に対する $t=PAv$ の余核像は
$H^0(C,K_C\otimes\eta)$ の標準的な一次元部分空間を定める。
十八点除数は直接
\[
 I_{D_{\rm mov}}=(\det Q,\operatorname{adj}(Q)t)
\]
で復元される。
その最小分解の次数は、生成元が $(5,5,5,6)$、syzygyが $(7,7,7)$。
三つの五次生成元は固定三角形の三頂点で可逆な対角評価行列を与える。
従って $h^0(I_D(5))=3$、$h^0(I_Z(5))=0$。
固定三点は核pencilだけから、移動十八点はこの標準切断から得られる。

補助六次曲線は、標準形で $v=e_0$ と正規化し、$s=v^tAv$ と置くと
\[
 \det A\det(RA^{-1}R^t)
 =s\det Q-t^t\operatorname{adj}(Q)t
\]
から得られる。この零曲線は $v$ の非零定数倍に依存せず、$Q$ と標準直線から計算できる。

\paragraph{次に調べる価値のある問題.}
十次元の余因子成分と三次元トーラスの表示により、有理性・不変量・分岐の計算が具体化した。
当面は九余因子方程式の大域的な消去や低次の不変量を調べることができるが、
形式局所解を有理解と取り違えないことが必要である。
全monodromyと反対合の標識忘却には、標準形が与える固定三角形以外の情報が要る。
また、九条件を単なる階数七条件から自動的に導いてはならない。
階数七・正則pencilの三十三次元軌跡と、本研究の二十四次元枠付き像は異なる。
